Kevin Kelleher suggested an interesting way to compare programming
languages: to describe each in terms of the problem it fixes. The
surprising thing is how many, and how well, languages can be described
this way.

\textbf{Algol:} Assembly language is too low-level.

\textbf{Pascal:} Algol doesn't have enough data types.

\textbf{Modula:} Pascal is too wimpy for systems programming.

\textbf{Simula:} Algol isn't good enough at simulations.

\textbf{Smalltalk:} Not everything in Simula is an object.

\textbf{Fortran:} Assembly language is too low-level.

\textbf{Cobol:} Fortran is scary.

\textbf{PL/1:} Fortran doesn't have enough data types.

\textbf{Ada:} Every existing language is missing something.

\textbf{Basic:} Fortran is scary.

\textbf{APL:} Fortran isn't good enough at manipulating arrays.

\textbf{J:} APL requires its own character set.

\textbf{C:} Assembly language is too low-level.

\textbf{C++:} C is too low-level.

\textbf{Java:} C++ is a kludge. And Microsoft is going to crush us.

\textbf{C\#:} Java is controlled by Sun.

\textbf{Lisp:} Turing Machines are an awkward way to describe
computation.

\textbf{Scheme:} MacLisp is a kludge.

\textbf{T:} Scheme has no libraries.

\textbf{Common Lisp:} There are too many dialects of Lisp.

\textbf{Dylan:} Scheme has no libraries, and Lisp syntax is scary.

\textbf{Perl:} Shell scripts/awk/sed are not enough like programming
languages.

\textbf{Python:} Perl is a kludge.

\textbf{Ruby:} Perl is a kludge, and Lisp syntax is scary.

\textbf{Prolog:} Programming is not enough like logic.
